\documentclass[12pt]{article}
\usepackage{amsmath,amssymb,amsthm}
\usepackage{geometry}
\usepackage{enumitem}
\usepackage{microtype}
\usepackage{booktabs}

%% hyperref MUST come last; hidelinks removes all red boxes
\usepackage[hidelinks]{hyperref}

\geometry{margin=1.25in}

\newtheorem{theorem}{Theorem}[section]
\newtheorem{lemma}[theorem]{Lemma}
\newtheorem{corollary}[theorem]{Corollary}
\newtheorem{definition}{Definition}[section]
\newtheorem{assumption}{Assumption}[section]
\newtheorem{invariant}{Invariant}[section]
\newtheorem{normalform}{Normal Form}[section]
\theoremstyle{remark}
\newtheorem{remark}{Remark}[section]

\title{\textbf{Principia Orthogona}\\[0.5em]
\large Volume One: The Mathematics of Generative Transitions}

%% \thanks must be INSIDE \author{}, not on a separate line
\author{Pablo Nogueira Grossi\thanks{%
  Academic Coordinator, UCEDA School (ESL), Elizabeth, NJ\@.
  Newark, NJ, USA\@.
  Email: \texttt{pablogrossi@hotmail.com}}}

\date{2026}

\begin{document}

\maketitle

\begin{center}
\emph{Dedicated to my children Vic~(R.I.P.), Giulia,
  Alice~(R.I.P.), Sarah~(R.I.P.), and David.}\\[4pt]
\emph{Once tiny, always strong.}
\end{center}

\vspace{1.5em}

\tableofcontents
\newpage

%-----------------------------------------------------------------------
\section*{Preface}
%-----------------------------------------------------------------------

This volume develops a unified mathematical framework for generative
transitions: localized geometric events in which a trajectory undergoes
compression, curvature intensification, loss of injectivity, and
stabilization.

The central object is the operator sequence
\[
C \;\rightarrow\; K \;\rightarrow\; F \;\rightarrow\; U,
\]
acting on trajectories in a Riemannian manifold.

Several results are presented with proof sketches rather than complete
proofs. These include the symplectic preservation across the fold, the
classification theorem for admissible singularities, the existence and
uniqueness of the unfolding map, and the stability of $\kappa^*$ under
perturbations. Each can be completed with standard tools from
differential geometry, singularity theory, and variational analysis;
full proofs lie beyond the scope of this foundational volume.

This volume does not claim that all physical, biological, cognitive,
or economic systems instantiate this structure. It provides the
mathematical substrate on which such questions can be posed rigorously.
Cross-domain instantiation is the subject of Volume Three.

This volume is the first in a three-part series. The subsequent
papers---\emph{Generative Contact Mechanics: A Geometric Framework for
Dissipative Systems with Structured Limit Cycles} and \emph{The dm3
Operator: Explicit Toy Model and Global Dynamical Analysis}---develop
the full contact-geometric realization of the generative transition
theory introduced here.

%-----------------------------------------------------------------------
\section{Overview}
%-----------------------------------------------------------------------

Let $(X, g)$ be a smooth Riemannian manifold. A \emph{generative
transition} is a localized event along a trajectory
$\gamma : [0,T] \to X$ consisting of four sequential phases:
compression, curvature induction, folding, and stabilization.

The theory is organized around the composite operator
\[
\mathcal{G} = U \circ F \circ K \circ C : X \to X.
\]

The manifold $N = (M, \Phi, F, K, g)$ serves as the canonical
instance, where $M$ is the geometric substrate, $\Phi$ the stability
functional, $F$ the directional field, $K$ the constraint operator,
and $g$ the generative operator.

The generative flow on $M$ is governed by
\[
\frac{dx}{dt} = -\nabla \Phi(x) + F(x) + \eta(t),
\]
where $\eta(t)$ is a stochastic perturbation.

%-----------------------------------------------------------------------
\section{Minimal Mathematical Assumptions}
%-----------------------------------------------------------------------

The operator sequence is defined under the weakest assumptions
necessary for the constructions of Section~\ref{sec:operators} to be
well-posed.

\begin{assumption}[Manifold Structure]\label{ass:manifold}
The state space $X$ is a smooth, finite-dimensional Riemannian
manifold, locally compact and second-countable.
\end{assumption}

\begin{assumption}[Trajectory Regularity]\label{ass:traj}
A system trajectory $\gamma : [0,T] \to X$ is piecewise $C^2$,
locally non-degenerate ($\|\dot\gamma(t)\| \neq 0$ a.e.), and has
bounded curvature on compact intervals prior to folding events.
\end{assumption}

\begin{assumption}[Compression Feasibility]\label{ass:compress}
There exists a lower-dimensional submanifold $X_C \subset X$ and a
Lipschitz continuous projection $C : X \to X_C$ satisfying the
non-collapse condition
\[
d\bigl(C(x_1), C(x_2)\bigr) \;\ge\; \delta\, d(x_1, x_2)
\]
for some $\delta > 0$ and all $x_1, x_2$ in a compact neighborhood.
\end{assumption}

\begin{remark}
Assumption~\ref{ass:compress} is a bi-Lipschitz embedding condition
on the compression map, ensuring distinct trajectories remain
distinguishable after compression. This is stronger than mere Lipschitz
continuity and should be understood as a lower bound on distance
preservation.
\end{remark}

\begin{assumption}[Curvature Threshold]\label{ass:kappa}
The critical curvature is defined intrinsically by the focal radius:
\[
\kappa^*(x) \;=\; \frac{1}{\operatorname{foc}(x)},
\]
where $\operatorname{foc}(x)$ is the distance from $x$ to the first
focal point along the normal direction. In the presence of positive
sectional curvature $K_{\mathrm{sec}}(x) > 0$, this is modified by
the Rauch comparison theorem to
\[
\kappa^*(x) \;=\; \min\!\left(\|II_x\|,\;
  \sqrt{K_{\mathrm{sec}}(x)}\right),
\]
where $\|II_x\|$ is the norm of the second fundamental form.
For $K_{\mathrm{sec}}(x) \le 0$, we have $\kappa^*(x) = \|II_x\|$.
\end{assumption}

\begin{assumption}[Folding Well-Posedness]\label{ass:fold}
The folding operator $F : X_C \to X_F$ satisfies: the Jacobian $dF$
loses rank by exactly $1$ at fold points; the fold is local; and the
fold produces a finite number of branches.
\end{assumption}

\begin{assumption}[Morse Stability Functional]\label{ass:morse}
The stability functional $\Phi : X \to \mathbb{R}$ is $C^2$, bounded
below on compact subsets, and \emph{Morse}: all critical points are
non-degenerate, i.e., $\nabla^2\Phi(x^*) \succ 0$ at every local
minimum~$x^*$.
\end{assumption}

%-----------------------------------------------------------------------
\section{Operator Definitions}\label{sec:operators}
%-----------------------------------------------------------------------

\subsection{State Space}

Let $X$ be as in Assumption~\ref{ass:manifold}. A system trajectory
is a piecewise $C^2$ map $\gamma : [0,T] \to X$. The goal is to
describe local transitions in $\gamma$ via the sequence
$C \to K \to F \to U$.

\subsection{Compression Operator \texorpdfstring{$C$}{C}}

\begin{definition}[Compression]
A \emph{compression operator} is a map $C : X \to X_C$ with
$\dim(X_C) < \dim(X)$ satisfying Assumption~\ref{ass:compress}.
\end{definition}

$C$ reduces degrees of freedom while retaining essential local
structure.

\subsection{Curvature Operator \texorpdfstring{$K$}{K}}

\begin{definition}[Curvature Operator]
Given a compressed trajectory $\gamma_C : [0,T] \to X_C$, the
curvature operator $K : X_C \to X_C$ modifies the tangent field by
\[
\frac{d}{ds}\bigl(K(\gamma_C)(s)\bigr)
\;=\; \dot\gamma_C(s) + \alpha(s)\,n(s),
\]
where $n(s)$ is the unit normal and
\[
\alpha(s) \;=\; \lambda\,\bigl(\kappa^*(\gamma_C(s))
  - \kappa(s)\bigr)_+, \quad \lambda > 0.
\]
\end{definition}

$K$ is a curvature-inducing flow that drives curvature monotonically
toward $\kappa^*$ but never beyond it. Consequently, $K$ alone does
not trigger folding; the sequence is a \emph{hybrid dynamical system}
with a switching condition at $\kappa^*$.

\subsection{Folding Operator \texorpdfstring{$F$}{F}}

\begin{definition}[Folding Map]
Let $\gamma_K = K(\gamma_C)$. A \emph{fold} occurs at $s_0$ when
$|\kappa_K(s_0)| = \kappa^*(\gamma_K(s_0))$. The folding operator
$F : X_C \to X_F$ acts by
\[
F(\gamma_K(s)) \;=\; \gamma_K(s) - \beta(s)\,n(s),
\]
where
\[
\beta(s) \;=\;
\begin{cases}
0, & |\kappa_K(s)| < \kappa^*, \\
\mu\bigl(|\kappa_K(s)| - \kappa^*\bigr), & |\kappa_K(s)| \ge \kappa^*.
\end{cases}
\]
Here $\mu > 0$ controls fold depth.
\end{definition}

At a fold point $s_0$,
$\operatorname{rank}(dF_{\gamma_K(s_0)}) = \dim(X_C) - 1$.

\subsection{Unfolding Operator \texorpdfstring{$U$}{U}}

\begin{definition}[Unfolding Map]
The unfolding operator $U : X_F \to X$ is
\[
U(x_F) \;=\; \operatorname*{arg\,min}_{y \in \mathcal{N}(x_F)}
  \Phi(y),
\]
where $\mathcal{N}(x_F)$ is a sufficiently small neighborhood of
$x_F$ and $\Phi$ satisfies Assumption~\ref{ass:morse}.
\end{definition}

The unfolding process is realised by the gradient flow
$\dot y = -\nabla\Phi(y)$, $y(0) = x_F$, which converges to a
non-degenerate local minimum~$x^*$.

\subsection{Compatibility of \texorpdfstring{$K$}{K} and
\texorpdfstring{$F$}{F}}

\begin{theorem}[Sequential Consistency]
If $K$ is applied until curvature reaches $\kappa^*$, then $F$ is
well-defined, produces a finite branch set, and induces a
rank-deficient Jacobian at fold points.
\end{theorem}

\begin{proof}[Proof sketch]
$K$ drives curvature monotonically to $\kappa^*$ via the gain field
$\alpha(s)$. At $\kappa^*$, the term $\beta(s)$ becomes nonzero,
introducing local non-injectivity. The normal direction collapses,
reducing Jacobian rank by~$1$. Locality and the finite critical-point
structure of $\kappa^*$ (via the Morse assumption on $\Phi$) imply
finitely many branches.
\end{proof}

%-----------------------------------------------------------------------
\section{Falsifiability Conditions}
%-----------------------------------------------------------------------

The framework is falsifiable at four levels.

\paragraph{Compression.}
If empirical data show expansion under $C$, or collapse of distinct
trajectories, the model fails.

\paragraph{Curvature.}
If a fold occurs at curvature strictly below $\kappa^*$, or curvature
exceeds $\kappa^*$ without folding, the model is falsified. ($\kappa^*$
is computed independently via the focal radius, not post-hoc from fold
observations.)

\paragraph{Folding.}
If empirical fold events do not correspond to Jacobian rank loss, or
produce infinitely many branches, the model fails.

\paragraph{Sequence order.}
If transitions occur in a different order, or stabilization occurs
without folding, the model is falsified.

%-----------------------------------------------------------------------
\section{Structural Theorems}
%-----------------------------------------------------------------------

\begin{theorem}[Existence and Well-Posedness]
Under Assumptions~\ref{ass:manifold}--\ref{ass:morse}, the composite
operator $\mathcal{G} = U \circ F \circ K \circ C$ is well-defined on
any piecewise $C^2$ trajectory.
\end{theorem}

\begin{theorem}[Local Determination]
The action of $\mathcal{G}$ on $\gamma$ is determined entirely by the
local geometry of $\gamma$ in a neighborhood of the fold point.
\end{theorem}

\begin{theorem}[Non-Commutativity]
The operators $C$, $K$, $F$, $U$ do not commute; the sequence is
order-dependent.
\end{theorem}

\begin{theorem}[Irreducibility]
No operator in the sequence $C \to K \to F \to U$ can be removed
without altering the qualitative structure of the transition.
\end{theorem}

\begin{theorem}[Finite Branching]
The branch set $B = \{F(\gamma_K(s_i)) :
|\kappa_K(s_i)| = \kappa^*(\gamma_K(s_i))\}$ is finite.
\end{theorem}

\begin{remark}
Finiteness follows from the Morse condition on $\Phi$
(Assumption~\ref{ass:morse}), which prevents fourth-order degeneracies
and ensures isolated critical points, together with the transversality
of $\gamma$ to the fold locus (a generic condition to be assumed
explicitly in applications).
\end{remark}

%-----------------------------------------------------------------------
\section{Canonical Examples}
%-----------------------------------------------------------------------

\paragraph{Planar curve with curvature-driven flow.}
The simplest nontrivial system: $\gamma : [0,T] \to \mathbb{R}^2$,
with $C$ an orthogonal projection, $K$ the curvature-inducing flow,
$F$ activated at $\kappa^*$, and $U$ gradient descent on a local
$\Phi$.

\paragraph{Elastic rod under compression.}
An elastic rod minimising Euler--Bernoulli energy
$E[\gamma] = \int \kappa(s)^2\,ds$. Axial loading provides $C$;
buckling provides $K$; the critical load is $\kappa^*$; post-buckling
configuration provides $U$. This is the cleanest physical realisation
of the curvature threshold.

\paragraph{Gradient flow on a double-well potential.}
$\Phi(x) = (x^2 - 1)^2$. The unstable equilibrium at $x = 0$ is the
fold point; gradient flow selects one of $x = \pm 1$. Illustrates
finite branching and stability selection.

\paragraph{Saddle-node bifurcation.}
$\dot x = \mu - x^2$, $\dot y = -y$. Projection onto the slow
manifold provides $C$; approach to the fold point provides $K$; loss
of the slow manifold at $\mu = 0$ provides $F$; flow to the stable
branch provides $U$.

\paragraph{Plasma sheet and market manifolds (Volume Three).}
Magnetic field-line reconnection and price-trajectory bifurcations are
conjectured instantiations. Each requires domain-specific construction
of $(X, g, \Phi)$ and empirical computation of $\kappa^*$. These are
deferred to Volume Three.

%-----------------------------------------------------------------------
\section{Analytical Invariants}
%-----------------------------------------------------------------------

\begin{invariant}[Ambient Dimension]
$\dim(X)$ is preserved by $\mathcal{G}$. No claim is made about the
dimension of the image of a specific trajectory.
\end{invariant}

\begin{invariant}[Codimension of the Fold]
$\operatorname{codim}(F(\gamma)) = 1$, following from the rank-loss
condition.
\end{invariant}

\begin{invariant}[Critical Curvature Threshold]
$\kappa^*(x)$ is a geometric property of the manifold, invariant under
the operator sequence.
\end{invariant}

\begin{invariant}[Curvature Sign]
$\operatorname{sgn}(\kappa(s))$ is preserved under $K$ and $F$; only
magnitude changes.
\end{invariant}

\begin{invariant}[Injectivity Before Threshold]
For all $s$ with $|\kappa(s)| < \kappa^*(s)$, the trajectory remains
injective under $C$, $K$, $F$.
\end{invariant}

\begin{invariant}[Rank Deficiency at Fold]
$\operatorname{rank}(dF) = \dim(X_C) - 1$ at every fold point.
\end{invariant}

\begin{invariant}[Energy Monotonicity under $U$]
$\Phi(U(x)) \le \Phi(x)$, with strict inequality unless $x$ is
already a local minimum.
\end{invariant}

\begin{remark}
An energy-gap invariant $\Delta\Phi = \Phi(x_2) - \Phi(x_1)$ is
preserved under $C$, $K$, $F$ only if $\Phi$ is constant along the
fibers of each operator. This is not guaranteed in general and is not
claimed as a universal invariant.
\end{remark}

%-----------------------------------------------------------------------
\section{Normal Forms}
%-----------------------------------------------------------------------

Two transitions $\mathcal{G}_1$, $\mathcal{G}_2$ are equivalent if
there exists a diffeomorphism $\psi : X \to X$ with
$\mathcal{G}_2 = \psi \circ \mathcal{G}_1 \circ \psi^{-1}$.

\begin{normalform}[Compression]
$C_{\mathrm{NF}} = \pi_k : \mathbb{R}^n \to \mathbb{R}^k$, the
standard coordinate projection.
\end{normalform}

\begin{normalform}[Curvature]
$K_{\mathrm{NF}}(s) = (s,\, \alpha s^2)$, $\alpha > 0$. This
produces curvature
$\kappa(s) = 2\alpha(1 + 4\alpha^2 s^2)^{-3/2}$, which approximates
$2\alpha$ near $s = 0$.
\end{normalform}

\begin{normalform}[Folding --- Whitney Fold]
$F_{\mathrm{NF}}(u,v) = (u, v^2)$. This is the unique (up to
diffeomorphism) normal form for a codimension-1 fold singularity.
\end{normalform}

\begin{normalform}[Unfolding]
$U_{\mathrm{NF}}(x) = 0$, the canonical stabilization map arising
from $\Phi_{\mathrm{NF}}(x) = x^2$.
\end{normalform}

%-----------------------------------------------------------------------
\section{Singularity Classification}
%-----------------------------------------------------------------------

\subsection{Admissible Singularities}

Given the constraints of rank-1 loss, finite branching, and the Morse
condition on $\Phi$, the admissible singularities are precisely:

\begin{enumerate}
\item Fold $A_1$: $(u,v) \mapsto (u, v^2)$
\item Cusp $A_2$: $(u,v) \mapsto (u, v^3 + uv)$
\item Swallowtail $A_3$: $(u,v) \mapsto (u, v^4 + uv^2 + \beta v)$
\end{enumerate}

Higher $A_k$ ($k \ge 4$) require a $k$-parameter unfolding. The Morse
condition on $\Phi$ restricts the unfolding to at most three
parameters, excluding $A_4$ and above.

\subsection{Classification Conditions}

Let $\Delta(s) = \kappa(s) - \kappa^*(s)$.

\begin{center}
\begin{tabular}{lll}
\toprule
Type & Conditions & Normal form \\
\midrule
$A_1$ (fold) & $\Delta(s_0)=0$, $\Delta'(s_0)\ne 0$ & $(u,v^2)$ \\
$A_2$ (cusp) & $\Delta=\Delta'=0$, $\Delta''\ne 0$ & $(u,v^3+uv)$ \\
$A_3$ (swallowtail) & $\Delta=\Delta'=\Delta''=0$, $\Delta'''\ne 0$
  & $(u,v^4+uv^2+\beta v)$ \\
\bottomrule
\end{tabular}
\end{center}

\begin{theorem}[Classification of Generative Transitions]
Every admissible generative transition $\mathcal{G} = U \circ F \circ
K \circ C$ is $\mathcal{A}$-equivalent to exactly one of $A_1$,
$A_2$, $A_3$.
\end{theorem}

\begin{proof}[Proof sketch]
Rank loss is exactly $1$, restricting to the $A_k$ series. The Morse
condition on $\Phi$ limits the unfolding to at most three parameters.
Transversality of $\gamma$ to the fold locus (generic assumption)
ensures isolated fold points. Thus $k \le 3$.
\end{proof}

\subsection{Bifurcation Stratification}

Parameter space $\Theta \subset \mathbb{R}^p$, $p \le 3$, decomposes
into:
\begin{itemize}
\item $\Theta_1$: generic fold region (codimension 0 in $\Theta$),
\item $\Theta_2$: cusp stratum (codimension 1),
\item $\Theta_3$: swallowtail stratum
  (codimension 2; a single point when $p=3$).
\end{itemize}

%-----------------------------------------------------------------------
\section{Metric Geometry of \texorpdfstring{$\kappa^*$}{kappa*}}
\label{sec:metric}
%-----------------------------------------------------------------------

\subsection{Intrinsic Definition}

\[
\kappa^*(x) \;=\; \frac{1}{\operatorname{foc}(x)},
\]
where $\operatorname{foc}(x)$ is the focal radius. For a curve
embedded in $(X,g)$, this equals $\|II_x\|$ (the norm of the second
fundamental form) when $K_{\mathrm{sec}} \le 0$.

\subsection{Sectional Curvature Correction}

By the Rauch comparison theorem, for variable positive sectional
curvature:
\[
\kappa^*(x) \;=\; \min\!\left(\|II_x\|,\;
  \sqrt{K_{\mathrm{sec}}(x)}\right).
\]
Positive ambient curvature lowers the threshold for folding.

\subsection{Computability}

Given $\gamma$ in $(X,g)$:
\begin{enumerate}
\item Compute $\|II_x\|$ from the second fundamental form.
\item Compute $K_{\mathrm{sec}}(x)$ from the curvature tensor.
\item Apply the formula above.
\end{enumerate}

$\kappa^*$ is stable:
$\delta\kappa^* = O(\delta g) + O(\delta II)
  + O(\delta K_{\mathrm{sec}})$.

%-----------------------------------------------------------------------
\section{Variational Principles}
%-----------------------------------------------------------------------

Define the action functional
\[
\mathcal{S}[\gamma] \;=\; \int_0^T
\Bigl(
\tfrac{1}{2}\|P_\perp \dot\gamma\|^2
+ \tfrac{\lambda}{2}(\kappa^* - \kappa)_+^2
+ \mu\,\delta(|\kappa| - \kappa^*)
+ \Phi(\gamma)
\Bigr)\,ds.
\]

Each term corresponds to one operator:
\begin{itemize}
\item $L_C = \tfrac{1}{2}\|P_\perp\dot\gamma\|^2$: minimise
  orthogonal kinetic energy (compression).
\item $L_K = \tfrac{\lambda}{2}(\kappa^*-\kappa)_+^2$: minimise
  curvature deficit.
\item $L_F = \mu\,\delta(|\kappa|-\kappa^*)$: singular activation
  at threshold.
\item $L_U = \Phi(\gamma)$: minimise stability potential.
\end{itemize}

The delta term places the framework in the class of
\emph{free-discontinuity variational problems}
(cf.\ Ambrosio--Tortorelli, Mumford--Shah, Francfort--Marigo).
It produces the jump condition
\[
\left[\frac{\partial L}{\partial \dot\gamma}\right]_{s_0^-}^{s_0^+}
= \mu\,n(s_0),
\]
which is the variational counterpart of the fold map.

The full generative transition is therefore a
\emph{piecewise-smooth extremal} of $\mathcal{S}$.

%-----------------------------------------------------------------------
\section{Hamiltonian Discontinuities and Symplectic Geometry}
%-----------------------------------------------------------------------

\subsection{Phase Space}

The phase space is $T^*X$ with canonical coordinates $(\gamma, p)$
and symplectic form $\omega = d\gamma \wedge dp$.

\subsection{Smooth Flow}

For $s \ne s_0$, the system obeys Hamilton's equations and the flow
is symplectic: $\mathcal{F}_t^*\omega = \omega$.

\subsection{Momentum Jump at the Fold}

The delta Lagrangian produces the impulse
\[
p(s_0^+) - p(s_0^-) \;=\; \mu\,n(s_0).
\]
Configuration $\gamma$ is continuous; momentum $p$ has a jump.

\subsection{The Fold as a Symplectic Map}

Define the fold map in phase space:
$\mathcal{F} : (\gamma, p) \mapsto (\gamma, p + \mu n)$.

\begin{theorem}[Symplectic Preservation]
$\mathcal{F}^*\omega = \omega$.
\end{theorem}

\begin{proof}
$d\gamma \wedge d(p + \mu n) = d\gamma \wedge dp
  + \mu\,d\gamma \wedge dn$.
Since $n$ depends only on $\gamma$, the term
$d\gamma \wedge dn = 0$.
Hence $\mathcal{F}^*\omega = \omega$.
\end{proof}

\subsection{Canonical Transformation}

The fold is generated by the distributional function
$S(\gamma) = \mu\,\Theta(|\kappa(\gamma)| - \kappa^*)$, so that
$p^+ = p^- + \partial S/\partial\gamma$.

\subsection{Full Transition}

Let $\Phi_t$ be the pre-fold Hamiltonian flow, $\mathcal{F}$ the fold
map, $\Psi_t$ the post-fold flow. The consistency condition from
Section~4 (basin structure) ensures the output of $\mathcal{F}$ lies
in the domain of $\Psi_t$. Then
\[
\mathcal{H} = \Psi_t \circ \mathcal{F} \circ \Phi_t
\]
is a piecewise-smooth symplectic map: $\mathcal{H}^*\omega = \omega$.

\subsection{Singularity Type and Momentum Structure}

\begin{itemize}
\item $A_1$ fold: single momentum jump.
\item $A_2$ cusp: momentum jump with first-order tangency constraint.
\item $A_3$ swallowtail: momentum jump with second-order tangency.
\end{itemize}

%-----------------------------------------------------------------------
\section{Connection to Generative Contact Mechanics}
\label{sec:connection}
%-----------------------------------------------------------------------

This section makes precise the relationship between the
singularity-theoretic framework developed in this volume and the
contact-geometric framework of \emph{Generative Contact Mechanics}
(Paper~1) and \emph{The dm3 Operator: Explicit Toy Model}
(Paper~2). We show that the operator sequence
$C \to K \to F \to U$
is the geometric precursor of the dm3 operator algebra, and that the
curvature threshold $\kappa^*$ corresponds canonically to the
embodiment threshold $\tau$.

\subsection{From Trajectories to Limit Cycles}

In this volume the central object is a trajectory
$\gamma : [0,T] \to (X,g)$ undergoing a localized generative
transition driven by curvature, injectivity loss, and stabilization.

In GCM the central object is a hyperbolic limit cycle
$\Gamma \subset S$ embedded in a dissipative dynamical system
satisfying Axioms~1--8 of the dm3 framework.

The conceptual shift is:
\begin{itemize}[leftmargin=*]
  \item \emph{Principia Orthogona}: analyzes single transitions
        along trajectories.
  \item \emph{GCM}: studies persistent post-transition attractors
        and their algebraic interactions.
\end{itemize}

The mathematical link is that every admissible generative transition
induces a locally attracting invariant set, and conversely, every
dm3 limit cycle admits a neighborhood whose formation is governed by
a fold-type transition.

\subsection{Curvature Threshold vs.\ Embodiment Threshold}

In this volume, folding is triggered when curvature reaches the
intrinsic geometric threshold
\[
  \kappa^*(x) = \frac{1}{\operatorname{foc}(x)},
\]
corrected by sectional curvature via the Rauch comparison theorem
(Section~\ref{sec:metric}).

In GCM, stochastic stability is governed by the embodiment threshold
\[
  \tau = \sqrt{c/\kappa},
\]
defined by the stochastic Lyapunov condition
$\mathcal{L}V \le -cV + \kappa\|\sigma\|^2$.

\begin{proposition}[Threshold Correspondence]\label{prop:threshold}
Let a generative transition produce a post-fold stabilized orbit
$\Gamma$ with Lyapunov function $V = d_g(\cdot,\Gamma)^2$. Then
the geometric condition $\kappa \uparrow \kappa^*$ is equivalent to
the dynamical condition that transverse contraction becomes dominant
over noise, yielding a finite embodiment threshold $\tau$.
\end{proposition}

\begin{proof}[Proof sketch]
As $\kappa \to \kappa^*$, the folding operator $F$ introduces
rank-1 loss and the post-fold unfolding $U$ selects a stable branch.
The resulting orbit $\Gamma$ has Floquet exponent
$\mu_{\max} < 0$ (transverse contraction). By the stochastic
Lyapunov condition of GCM, $\tau = \sqrt{c/\kappa}$ is finite
whenever $\mu_{\max} < 0$, i.e., precisely when curvature has
reached $\kappa^*$ and the fold has been completed. Thus
$\kappa^*$ is the geometric precursor of $\tau$: curvature
accumulation creates the conditions under which stochastic stability
becomes meaningful.
\end{proof}

\subsection{Fold Map vs.\ dm3 Operator}

The folding operator $F$ in this volume is a rank-1 singular map:
$\operatorname{rank}(dF) = \dim(X_C) - 1$, producing a finite
branch set classified by $A_1$--$A_3$ singularities.

In GCM, the dm3 operator
$\mathcal{A}_{\mathrm{dm3}} = \phi^{T^*/4}$
is the time-$(T^*/4)$ map of a smooth flow on the contact manifold
$M = S \times \mathbb{R}$, acting discretely along a limit cycle.

\begin{proposition}[Fold--dm3 Correspondence]\label{prop:fold}
The fold map $F$ is the piecewise-smooth, pre-contact limit of the
dm3 operator in the following sense.
\begin{enumerate}[label=\textup{(\roman*)}]
  \item $F$ enforces non-injectivity geometrically via rank loss;
        $\mathcal{A}_{\mathrm{dm3}}$ enforces recurrence dynamically
        via orbital stability.
  \item Under the contact extension $M = X \times \mathbb{R}$
        \textup{(}the construction of GCM, \S6\textup{)}, the
        impulsive momentum jump at the fold
        $p(s_0^+) - p(s_0^-) = \mu\,n(s_0)$
        corresponds to the contact Hamiltonian correction
        $H_{\mathrm{diss}} = -\gamma V e^{-\beta z}$,
        which drives trajectories toward $\Gamma$ off the orbit.
\end{enumerate}
\end{proposition}

\begin{proof}[Proof sketch]
Both the momentum jump and $H_{\mathrm{diss}}$ are corrections
that vanish on the invariant set and are nonzero off it. The
momentum jump is a distributional first-order correction in the
Hamiltonian sense; $H_{\mathrm{diss}}$ is its smooth
contact-geometric regularization as the contact manifold structure
is imposed via $M = X \times \mathbb{R}$.
\end{proof}

\subsection{Operator Grammar Correspondence}

The operator sequence of this volume maps onto the reduced dm3
operator grammar as follows.

\begin{center}
\begin{tabular}{ll}
\toprule
\emph{Principia Orthogona} & \emph{GCM / dm3} \\
\midrule
Compression $C$ & Basin contraction via $g_1$ \\
Curvature flow $K$ & Lyapunov descent $\dot{V} \le -cV$ \\
Fold $F$ & Contact bifurcation (Hopf, SN, NS) \\
Unfolding $U$ & Gradient flow to $\Gamma_{\mathrm{syn}}$ via $U_3$ \\
\bottomrule
\end{tabular}
\end{center}

The full dm3 grammar $g \to L \to R \to U$ extends this sequence
by adding multi-orbit coherence, resonance detection, and categorical
unification, which are intentionally absent from the present volume.

\subsection{Singularity Classification vs.\ Bifurcation Diagram}

This volume proves that admissible generative transitions are
precisely the Whitney singularities $A_1$, $A_2$, $A_3$.
Paper~2 proves that dm3 systems undergo exactly four bifurcations:
contact Hopf, saddle-node of limit cycles, Neimark--Sacker, and
slow-fast crossover.

\begin{remark}[Classification Correspondence]
The $A_1$--$A_3$ singularities classify the local geometry of dm3
bifurcations under the projection $M = S \times \mathbb{R} \to S$.
The fold $A_1$ corresponds to the contact Hopf and saddle-node
bifurcations (local rank loss in the radial direction); the cusp
$A_2$ corresponds to the Neimark--Sacker (loss of a second
direction at resonance); the swallowtail $A_3$ to the slow-fast
crossover (third-order degeneracy in the $z$-direction). Higher
singularities are excluded in both frameworks: here by the Morse
condition on $\Phi$; in dm3 by contact normal form rigidity
(Theorem~C of Paper~1).
\end{remark}

\subsection{Summary}

The three papers in this series have the following relationship.
\begin{itemize}[leftmargin=*]
  \item \textbf{This volume}: explains why folds must occur and
        classifies their local geometry.
  \item \textbf{Paper~1 (GCM)}: explains what stable structures
        exist after they occur, in the contact-geometric setting.
  \item \textbf{Paper~2 (dm3)}: proves that the resulting dynamics
        are computable, stable, and structurally closed.
\end{itemize}

No claim is made that every dm3 system arises from a literal
curvature fold in a physical manifold. The precise claim is:
every dm3 system admits a local geometric model whose generative
event is $\mathcal{A}$-equivalent to a fold-type transition governed
by a curvature threshold $\kappa^*$. This is supported by
Propositions~\ref{prop:threshold} and~\ref{prop:fold}.

%-----------------------------------------------------------------------
\section*{Closing Statement}
%-----------------------------------------------------------------------

This volume has developed a unified mathematical framework for
generative transitions. The central results are:
\begin{itemize}
\item a geometric foundation based on curvature, injectivity radius,
  and focal geometry;
\item a critical curvature threshold $\kappa^*$ defined intrinsically
  by the focal radius and corrected by the Rauch comparison theorem;
\item constructive definitions of operators $C$, $K$, $F$, $U$;
\item a free-discontinuity variational principle;
\item a Hamiltonian structure with a distributional generator at the
  fold;
\item a symplectic discontinuity preserving the canonical form;
\item a singularity classification restricted to the Whitney
  $A_1$--$A_3$ hierarchy; and
\item structural invariants holding across all admissible transitions.
\end{itemize}

The framework is placed within established mathematical traditions:
comparison geometry, singularity theory, geometric flows, variational
mechanics, and impulsive Hamiltonian systems.

This volume is deliberately limited in scope. It provides the
mathematical substrate on which cross-domain questions can be posed
rigorously. Volume Two develops discrete variational integrators and
numerical realisation. Volume Three instantiates the framework in
specific domains: plasma reconnection, market volatility manifolds,
and neural embedding geometry.

%-----------------------------------------------------------------------
\section*{Acknowledgments}
%-----------------------------------------------------------------------

The author acknowledges with gratitude the foundational influence of
the teachings of Paramahamsa Nithyananda and the yogic scriptural
tradition from which this work derives. The mathematical formalization
presented here is understood by the author as a translation of
principles encountered through that tradition.

This volume is the first in a series. The dm3 operator framework
developed in the companion papers \emph{Generative Contact Mechanics}
(Paper~1) and \emph{The dm3 Operator: Explicit Toy Model} (Paper~2)
constitutes the full mathematical realization of the generative
transition theory introduced here.

\end{document}
